% Pursuing Stacks (cotes 134-2 to 134-8): archive cote 134-5, site page 34 (PDF page 35), typescript page 221.
% typescript. Draft compared: ps4.tex:1640-1703.
% AI-transcribed by Claude Opus 5.5 (claude-opus-5-5, adaptive thinking, effort high), Message Batches API, batch ps2 (msgbatch_01NVwr1Yuo8qJf5rykeGBXPK), 2026-11-12; reviewed against both supplied scan halves; see CHECK_DECISIONS.tsv.
\scanpage{134-5/34}
\mspage{221}
means that for any b in B, the category
\[ {}_{b\backslash}A \overset{\mathrm{def}}{=} Ax_B({}_{b\backslash}B) \;=\; \text{category of pairs } (x,p)\hand{,} \text{ with } x \text{ in } A \]
\[ \text{and } p : b \hand{\rightarrow} \over{\unsure{g}}{i}(x) \]
\margin{$i(x)$}
is aspheric. (NB In the case of the functor g from B to A, the corres-%
ponding categories ${}_{a\backslash}B$ \inserted{even} have \hand{\uline{\textcolor{black}{initial}}} \hand{\uline{\textcolor{black}{objects}}}.) This condition comes in
here rather formally, we'll see later though that it has a \xout{very} \inserted[typed]{quite} remarkab-%
le interpretation, in terms of a very strong property of cohomological
``cofinality'' of the functor i, implying the usual notion of i being a
``cofinal'' functor\hand{,}or A being ``cofinal''in B (namely i giving rise to
isomorphisms $\underset{\hand{\longrightarrow}}{\lim}_A \hand{\longrightarrow} \underset{\hand{\longrightarrow}}{\lim}_B$ for any \xout{indexed system} direct system
B $\hand{\longrightarrow}$M with values in a category M admitting direct limits)\hand{,} as the ``dime\hand{n}-%
sion zero'' shadow of this ``all dimensions'' property. To make an anology
which will acquire more precise meaning later, the asphericity property
for a map i in (Cat) can be viewed as a (slightly weakened) version of
a \hand{\uline{\textcolor{black}{proper map with aspheric fibers}}}, whereas the coasphericity property
appears as the corresponding version of a \hand{\uline{\textcolor{black}{smooth map with aspheric fibers}}}.
The qualification ``slightly weakened'' reflects notably in the fact that
the notions \xout{considered} \inserted[typed]{of asphericity and coasphericity} are \hand{\uline{\textcolor{black}{not}}} stable under arbitrary base change -
but rather\hand{,} asphericity for a map is stable under base change by \xout{a} \hand{\uline{\textcolor{black}{fibra-%
tion functors}}}(more generally, by \xout{a} smooth maps in (Cat)), whereas
coasphericity is stable under base change by \hand{\uline{\textcolor{black}{cofibration functors}}}
(more generally, by proper maps in (Cat)). Thus, whereas a functor
i: A $\hand{\longrightarrow}$ B which is an equivalence of categories is clearly both aspheric
and coaspheric, this property \xout{does} \inserted[typed]{is} not preserved by arbitrary base
change e.g. passage to fibers : \inserted{e.g.} some fibers may be empty, and therefore
are neither adpheric nor coaspheric !

The most comprehensive property \xout{including both} for a map in (Cat),
implying both asphericity and coasphericity, is to be in U$\underline{W}$ i.e. a
``\hand{\uline{\textcolor{black}{universal weak equivalence}}}'' - namely it is in $\underline{W}$ and remains so after
any base change. These maps deserve the name of ``\hand{\uline{\textcolor{black}{trivial Serre fibrati-%
ons}}}''. They include\inserted[typed]{all} maps with aspheric fibers which are either ``proper''
or ``smooth'', for instance those which are either \inserted[typed]{co}fibering or \xout{co}fibering
functors (in the usual sense of category theory). We'll come back upon
these notions in a \xout{more} systematic way in the next part of the reflecti\hand{\unsure{ons}}

74. As announced yesterday, I would like still to make an overdue
digression on products of test categories, before embarking on the
notion of asphericity structure. It will be useful to begin with some

\margin{34}
